\chapter{Laboratory Infrastructure, Financial Audit \& Recommendations}
\label{chap:infrastructure}

\section{Establishment of the Central DNAx Laboratory}
To house the physical synthesis, nanomaterial encapsulation, and forensic verification workflows, a specialized, high-security facility — \textbf{DNAx Laboratory (Room 309-B)} — was formally sanctioned and established at ASAB-NUST (Office Order No. 0958/02/ASAB/Accounts/1298, attached in Annexure~A).

\begin{figure}[H]
    \centering
    % Placeholder for 3D CAD Lab Floorplan (May Report Fig 6)
    \fbox{\parbox{0.85\textwidth}{\centering \vspace{1.8cm} \textbf{[Figure 8.1: 3D CAD Architectural Floor Plan and Equipment Layout of the DNAx Laboratory (Room 309-B)]} \vspace{1.8cm}}}
    \caption{Architectural layout and zoning of the specialized DNAx research facility.}
    \label{fig:lab_cad}
\end{figure}

\subsection{Current Status of Laboratory Assets and Instrumentation}
\begin{enumerate}
    \item \textbf{24/7 High-Definition Surveillance and Network Infrastructure}: A dedicated closed-circuit television (CCTV) system and isolated local network infrastructure were installed to maintain continuous monitoring of sensitive biological storage areas and computing servers.
    \item \textbf{Clean-Air Laminar Flow Workstation}: A certified vertical laminar flow cabinet was installed and calibrated to establish a Class 100 sterile micro-environment, preventing airborne DNA cross-contamination during sample prep.
    \item \textbf{Dedicated Biological Cold Storage}: Dual-temperature cold storage facilities ($-20^\circ\text{C}$ and $+4^\circ\text{C}$) were operationalized for long-term stabilization of synthesized DNA taggants, primers, probes, and enzymes.
    \item \textbf{Real-Time micPCR Machine Acquisition}: Procurement of a specialized, multi-channel mic-Polymerase Chain Reaction (micPCR) thermal cycler was finalized via AXI channels, providing the primary thermal cycling and optical signal acquisition hardware.
\end{enumerate}

\section{Financial Audit and Budget Utilization Statement}
Table~\ref{tab:budget_statement} presents the formal financial audit and budget utilization statement for the first installment as of \textbf{August 25, 2026} (Third Technical Panel Review).

\begin{table}[H]
\centering
\small
\caption{Comprehensive Budget Utilization Statement (1st Installment, as of 25 August 2026)}
\label{tab:budget_statement}
\begin{tabularx}{\textwidth}{l r r r r}
\toprule
\textbf{Budget Category} & \textbf{Approved Cost} & \textbf{1st Installment} & \textbf{Funds Utilized} & \textbf{Balance at NUST} \\
& \textbf{(PKR)} & \textbf{Received (PKR)} & \textbf{(PKR)} & \textbf{Account (PKR)} \\
\midrule
1. Equipment & 7,000,000 & 1,591,500 & 1,591,374 & 126 \\
2. Personnel (HR) & 5,280,000 & 3,240,000 & 3,240,000 & 0 \\
3. Lab Supplies \& Reagents & 3,500,000 & 2,393,500 & 2,390,579 & 2,921 \\
4. Field Operations & 402,000 & 0 & 0 & 0 \\
5. Data Management \& Analysis & 500,000 & 0 & 0 & 0 \\
6. Training \& Capacity Building & $\sim$400,000 & 0 & 0 & 0 \\
7. Monitoring \& Evaluation & $\sim$200,000 & 0 & 0 & 0 \\
8. Contingencies & $\sim$700,000 & 0 & 0 & 0 \\
9. NUST Institutional Overhead (15\%) & $\sim$2,997,000 & 1,275,000 & 1,275,000 & 0 \\
\midrule
\textbf{TOTAL} & \textbf{19,980,000} & \textbf{8,500,000} & \textbf{8,496,953} & \textbf{3,047} \\
\bottomrule
\end{tabularx}
\end{table}

As documented in Table~\ref{tab:budget_statement}, the first installment of \textbf{PKR 8,500,000} has been utilized with exceptional financial efficiency ($99.96\%$ utilization, with a remaining balance of only PKR 3,047), directly yielding the experimental breakthroughs documented in this report.

\section{Operational Bottlenecks and Strategic Recommendations}

\subsection{Identified Operational Bottlenecks}
\begin{enumerate}
    \item \textbf{Institutional Timings and Network Bandwidth}: Strict institutional working hours (9:00 AM to 5:00 PM) at ASAB-NUST restrict continuous multi-day synthesis runs, and limited local network bandwidth hinders bulk offline genomic database mirroring.
    \item \textbf{Coating and Uncoating Yield Repeatability}: While proof-of-concept encapsulation has been validated, further process optimization is needed to maximize recovery yields across complex aged substrates.
    \item \textbf{Explosive Testing Site Access}: Testing synthetic taggants under high-velocity detonation conditions requires remote military-grade explosive firing ranges not directly adjacent to the university campus.
    \item \textbf{Budget Depletion of First Installment}: Total funds from the first installment are fully exhausted. Without immediate capital injection, ongoing personnel retention and consumable procurements will be compromised.
\end{enumerate}

\subsection{Strategic Recommendations and Requests}
\begin{enumerate}
    \item \textbf{Release of Second Installment Funds}: The balance amount remaining with the sponsor (GHQ / AXI) is \textbf{PKR 3,380,000}. Immediate release of \textbf{PKR 2,080,000} is formally requested to maintain uninterrupted laboratory operations and research associate retention.
    \item \textbf{Pooling of Dedicated Coating Resources}: Additional capital allocation should be directed toward micro-extrusion and high-precision fluidic spray hardware for repeatable pilot-scale coating.
    \item \textbf{Specialized Working Authorizations}: NUST administration should be formally petitioned to grant 24/7 flexible laboratory access and dedicated high-speed fiber-optic bandwidth to the DNAx research team.
\end{enumerate}
